\subsection{Observable calibration and unknown initialization}
A covariance probe must be computable from the learner's observations.
Let its raw output be $Z=\boldsymbol1\nu+W_j$, where $\nu$ is the unknown target-induced level, $(C_T,W_j)$ is centered jointly Gaussian under regime $j$, and $\operatorname{Cov}_j(W_j)=\Sigma_j\succ0$.
Subtracting the unknown level $\nu$ would require unavailable target information.
Instead choose the first-difference matrix $D$ and use the observable contrast $Y=DZ=DW_j$.
The classical error-contrast construction~\cite{harville1974contrasts} gives covariance $K_j=D\Sigma_jD^\top$ and Gaussian affinity
\[
\rho_D=\det(K_0)^{1/4}\det(K_1)^{1/4}/\det[(K_0+K_1)/2]^{1/2}.
\]
Each directional error of its zero-threshold likelihood rule is at most $\rho_D$.
The common state at probe exit has conditional mean $m_j(Y)$ and explained variance $s_j^2=\operatorname{Var}_j(m_j(Y))$, computed from this projected law.
The supplement derives a linear-time two-filter implementation and a finite-dimensional information-cost bound for removing the target.
Every probe observation still consumes both budgets.

Let the initial error be an unknown fixed $e\in[-r,r]$, with public $r>0$.
In this subsection $j$ may range over any fixed finite collection of covariance models.
Fix a regime prior $\pi_j>0$, mixture law $g$, and model posterior $p_j(y)$.
Assume that (a) the probe and every continuing program are frozen before observing $Y$; (b) $e$ does not affect the contrast law; and (c) the continuing learners have deterministic linear weights and the common-state representation in Eq.~\eqref{eq:linear-program}.
Write the deterministic initialization contribution as $d_a(e)=\zeta_a e$ and let $k_{j,a}$ denote its common-state transfer coefficient under model $j$.
For each frozen post-probe program let $J_{j,a}^{r}=\zeta_a^2r^2+\operatorname{Var}_j(\sum_t\omega_{a,t}\bar X_t)$, including its remaining horizon.
Its even conditional score is
\begin{equation}
L_a(y)=\sum_jp_j(y)\{J_{j,a}^{r}+k_{j,a}^2[m_j(y)^2-s_j^2]\}.
\label{eq:conditional-even-score}
\end{equation}
This score uses neither the target nor the sign of the initial error.
Write $F_a(y)=\sum_jp_j(y)J_{j,a}^{r}$, let $u$ minimize $L_a$ and $c$ minimize $F_a$, and set $S_j(y)=m_j(y)^2-s_j^2$.
Both selectors use a fixed even tie rule.
For an even selector write $R_\pi(v)=\sup_{|e|\le r}R_{\pi,e}(v)$, and let $\varphi$ be the standard normal density.

\begin{theorem}[Conditional value and its state-information ceiling]
For a fixed paid probe and finite scalar continuing-program catalogue with the same deterministic learner weights under each regime,
\begin{equation}
\inf_v\sup_{|e|\le r}R_{\pi,e}(v)=\E_g\min_aL_a(Y),
\label{eq:conditional-minimax-value}
\end{equation}
where $v$ ranges over randomized $Y$-measurable selectors in the invariant-statistic information class.
The selector $u$ attains the value, and
\begin{align}
0\le R_\pi(c)-R_\pi(u)&\le K,\nonumber\\
K&=4\varphi(1)\sum_j\pi_j\operatorname{range}_a(k_{j,a}^2)s_j^2.
\label{eq:state-selection-ceiling-main}
\end{align}
For the running mean and a stable averaged affine-delay learner with affordable horizons proportional to budget scale $B$ and uniform persistence and impulse bounds, $K=O(B^{-2})$.
\label{thm:conditional-value}
\end{theorem}

Posterior-loss minimization is classical~\cite{moustakides2012joint,vasconcelos2020observation,reinhard2026minimax}; here endpoint averaging and evenness resolve the unknown initialization.
The supplementary proofs bound each regime's centered-square term separately, allowing different persistences, and give the vector covariance dual using classical minimax interchange~\cite{sion1958minimax}.
An explicit two-program running-mean construction attains a strictly positive $B^{-2}$ state-value limit, showing that this order is sharp.

For the best fixed risks $I_{\rm full},I_{\rm post}$ before and after payment and $R_{\rm rob}=\E_g\min_aL_a$, the paid gain decomposes as
\[
I_{\rm full}-R_{\rm rob}=V_{\rm reg}+V_{\rm state}-C-T,
\]
Here $V_{\rm reg}\ge0$ is nominal regime value, $V_{\rm state}\in[0,K]$ is extra state value, $C$ is the possibly signed remaining-horizon cost, and $|T|\le K/2$ corrects nominal posterior scores to actual risk.
The supplement proves this identity and the resulting paid-value bounds.

\subsection{When identification outperforms all fixed prefixes}
Theorem~\ref{thm:conditional-value} separates knowledge of the covariance regime from knowledge of the current noise state.
The latter has a second-order value for averaged learners; the former can change the leading risk coefficient.
For a running mean with $0\le\lambda\le1-\gamma<1$ and budgets $(\lfloor\beta_m B\rfloor,\lfloor\beta_e B\rfloor)$ for fixed $\beta_m,\beta_e>0$, define
\begin{align*}
a_a&=\min\{\beta_m/(h+q_a),\beta_e/b_a\},\\
K_j(a)&=\frac{q_a^{-1}+\theta_j(1+\lambda^{b_a})/(1-\lambda^{b_a})}{a_a}.
\end{align*}
Let $a_j^*$ minimize $K_j$, $K_\pi=\min_a\sum_j\pi_jK_j(a)$,
$\Delta=K_\pi-\sum_j\pi_jK_j(a_j^*)$, and
$A=\sum_j\pi_j[K_j(a_{1-j}^*)-K_j(a_j^*)]$.
The fixed comparison permits every affordable training prefix and stop, at the same learner and step.

\begin{corollary}[First-order separation and a no-gain boundary]
For a finite catalogue and bounded delay, a paid target-invariant probe with fixed directional error bound $0<\delta<1/2$, followed by commitment to $a_{\widehat j}^*$, certifies superiority to every fixed prefix and fixed random mixture for sufficiently large $B$ if $\Delta>\delta A$.
The certified gain is at least $(\Delta-\delta A)/B-C(n+1)/B^2$, where $n$ is probe length and $C$ is finite.
If the regimes share a leading minimizing action, no selector from the maximal-horizon post-probe catalogue plus stop following a length-$O(\log B)$ probe gains more than $O(\log B/B^2)$ over the best fixed prefix.
\label{cor:leading-prefix-separation}
\end{corollary}

The supplement proves the prefix expansion, gives explicit conservative remainder constants and a sufficient budget threshold, and supplies a rational witness on the public grid with a positive leading margin of approximately $0.07765$ at directional bound $0.025$.
Finite-budget admission uses the actual projected probe, exact horizons and signed risk bounds, rather than the asymptotic margin.
The Gaussian branch admits a probe only when its prior-weighted upper risk is below the best full-budget fixed-prefix risk; otherwise it executes that fixed prefix without probing.
The comparison concerns a declared learner and resource-program class, rather than unrestricted estimators.
Here a fixed prefix repeats one constant participation--spacing pair; an arbitrary time-varying nonadaptive schedule is a different comparator.

\subsection{Finite separation for averaging and TD critics}
Consider $Q=h=8$, initialization radius one, budgets $(16B,4B)$ and $q,b\in\{1,2,4,8\}$.
The covariance box is $\theta_0\in[0.04,0.06]$, $\theta_1\in[1.8,2.2]$, independently varying $\lambda_j\in[0.15,0.25]$, and prior $\pi\in[0.45,0.55]$.
Use eighteen paid unit-spacing observations from all actors, the target-invariant likelihood rule frozen at $(0.05,2,0.2)$, and commitments $(8,4)$ or $(2,2)$.
For each learner, let $I_{\rm open}$ be its best worst-initialization prior risk over every predetermined schedule, stopping prefix and independent random mixture, with actual covariance parameters available to that comparator.
For the averaged scalar TD(0) critic, use $e_t=\beta e_{t-1}+\eta\bar X_t$ with $\eta=0.1$, $\mu=1-\gamma=0.5$ and $\beta=1-\eta\mu$.

\begin{theorem}[Robust finite-budget separation]
Throughout the stated box, at $B=8192$ the paid rule satisfies
\[
R_{\rm paid}\le(1-\epsilon)I_{\rm open},
\]
with $\epsilon=0.0254$ for the running mean with eight withheld updates and $\epsilon=0.0186$ for the averaged TD critic with zero update delay.
All probe and scheduled-update costs are included.
\label{thm:finite-robust-separation}
\end{theorem}

The comparison allows arbitrary schedule length and order, including nonperiodic schedules.
Supplementary Section~\ref{sup-sec:td-extension} proves the result by direct Gaussian error bounds and a period-free resource envelope.
An explicit TD filter-boundary correction transfers the envelope beyond running means.

